% Part: proof-theory
% Chapter: sequent-calculus
% Section: translations

\documentclass[../../../include/open-logic-section]{subfiles}

\begin{document}

\olfileid[pt]{pt}{seq}{tr}

\olsection{Traduções entre \Log{G1c} e \Log{G3c}}

Mencionamos que tanto \Log{G1c} quanto \Log{G3c} são corretos e
completos para a lógica clássica, isto é, cada um demonstra exatamente
os sequentes verdadeiros em toda !!{structure}. Em particular, eles
demonstram os mesmos sequentes. Agora podemos demonstrar esse fato sem o
desvio pelos teoremas da correção e da completude, recorrendo apenas a
meios teórico-probatórios. Uma demonstração desse tipo fornece
\emph{transformações} que levam !!{proof}s em \Log{G1c} a !!{proof}s em
\Log{G3c} e vice-versa.

\begin{prop}
Se $\Log{G1c} \Proves S$, então $\Log{G3c} \Proves S$.
\end{prop}

\begin{proof}
  Mostramos que, se $\pi$ é !!a{proof} em \Log{G1c}, então existe
  !!a{proof}~$\pi'$ de~$S$ em \Log{G3c}. Fazemos isso por indução sobre
  $n = \pheight{\pi}$.

  Se $n = 0$, então $\pi$ é um axioma. O axioma $\lfalse \Sequent
  \quad$ também é um axioma em~\Log{G3c} (tome os contextos $\Gamma,
  \Delta$ como vazios). Em \Log{G3c}, são permitidos axiomas
  $!A \Sequent !A$ com $!A$ arbitrária, não apenas atômica. Mas
  demonstramos em \olref[ex]{prop:G3c-axioms} que todos esses sequentes
  são !!{provable} em~\Log{G3c}.
  
  Se $n>0$, $\pi$ termina em uma regra~$R$. A hipótese de indução
  garante que existem !!{proof}s das premissas dessa regra em~\Log{G3c},
  isto é, que as sub-!!{proof}s imediatas têm traduções em~\Log{G3c}. Se
  a regra~$R$ também é uma regra de \Log{G3c}, nós a aplicamos às
  !!{proof}s traduzidas das premissas. Quanto às regras que não são
  também regras de~\Log{G3c}, já mostramos que ao menos são admissíveis.
  Portanto, se a regra~$R$ era $\LeftR{\land}$ ou $\RightR{\lor}$,
  aplique \olref[adm]{prop:lor-G3-adm}. Se era uma regra de
  enfraquecimento, aplique \olref[adm]{prop:weak-G3c-adm}. Se era uma
  regra de contração, aplique \olref[inv]{prop:G3c-cont-adm}.
\end{proof}

\begin{prop}
Se $\Log{G3c} \Proves S$, então $\Log{G1c} \Proves S$.
\end{prop}

\begin{proof}
  Novamente, procedemos por indução sobre a !!{height} de
  !!a{proof}~$\pi$ de $S$. Todo axioma de \Log{G3c} pode ser !!{prove}d
  em~\Log{G1c} a partir de um axioma e de várias aplicações das regras
  \LeftR{\Weakening} e \RightR{\Weakening}:
  \[
  \Axiom$!C \fCenter !C$
  \RightLabel{\LeftR{\Weakening}}
  \Deduce$!C, \Gamma \fCenter !C$
  \RightLabel{\RightR{\Weakening}}
  \Deduce$!C, \Gamma \fCenter \Delta, !C$
  \DisplayProof
\qquad
  \Axiom$\lfalse \fCenter$
  \RightLabel{\LeftR{\Weakening}}
  \Deduce$\lfalse, \Gamma \fCenter$
  \RightLabel{\RightR{\Weakening}}
  \Deduce$\lfalse, \Gamma \fCenter \Delta$
  \DisplayProof
  \]
  Se $\pi$ termina em uma regra~$R$, a hipótese de indução fornece
  !!{proof}s das premissas em \Log{G1c}, e podemos aplicar a mesma
  regra~$R$ a essas !!{proof}s. As exceções são as regras
  \LeftR{\land} e \RightR{\lor}, que diferem entre \Log{G1c}
  e~\Log{G3c}. As versões de \Log{G3c} são deriváveis em \Log{G1c} da
  seguinte maneira:
  \[
  \Axiom$!A, !B, \Gamma \fCenter \Delta$
  \RightLabel{\LeftR{\land}}
  \UnaryInf$!A \land !B, !B, \Gamma \fCenter \Delta$
  \RightLabel{\LeftR{\land}}
  \UnaryInf$!A \land !B, !A \land !B, \Gamma \fCenter \Delta$
  \RightLabel{\LeftR{\Contraction}}
  \UnaryInf$!A \land !B, \Gamma \fCenter \Delta$
    \DisplayProof
\qquad
  \Axiom$\Gamma \fCenter \Delta, !A, !B$
  \RightLabel{\RightR{\lor}}
  \UnaryInf$\Gamma \fCenter \Delta, !A \lor !B, !B$
  \RightLabel{\RightR{\lor}}
  \UnaryInf$\Gamma \fCenter \Delta, , !A \lor !B, !A \lor !B$
  \RightLabel{\RightR{\Contraction}}
  \UnaryInf$\Gamma \fCenter \Delta, !A \lor !B$
    \DisplayProof
  \]
\end{proof}

\end{document}
